problem:  
Let $\{(M_i^n,g_i)\}$ be a sequence of closed Riemannian manifolds with $\mathrm{sec}_{g_i}\ge -1$ and $\operatorname{diam}(M_i,g_i)\le1$, collapsing in the Gromov–Hausdorff sense to a compact Alexandrov space $(X^k,d_X)$ with $k<n$.  
For each $x\in X$, fix points $p_i\in M_i$ converging to $x$, and consider the rescaled pointed universal covers $(\widetilde M_i, \varepsilon_i^{-2}g_i,p_i)$ with $\varepsilon_i\to0$, converging in the pointed Gromov–Hausdorff sense to $(Y_x,y_x)$, an Alexandrov space of nonnegative curvature.  

**Question:**  
Is it possible to assign canonically to each $(Y_x,y_x)$ a *locally compact metric group model* $(G_x,d_x,e_x)$ such that:

1. there exists $r>0$ and $\delta_i\to0$ with $B_{Y_x}(y_x,r)$ $\delta_i$‑Gromov–Hausdorff close to $B_{G_x}(e_x,r)$, where $G_x$ is equipped with a left‑invariant geodesic metric $d_x$;
2. the isometry type of $(G_x,d_x)$ depends only on the germ of $(Y_x,y_x)$—that is, it is canonically determined by the local geometry of the blow‑up limit;
3. as $x$ varies over $X$, the family $\{(G_x,e_x)\}_{x\in X}$ varies *continuously* in the pointed Gromov–Hausdorff topology (so nearby points of $X$ yield nearby metric group models).

Equivalently, does every collapsing sequence with a uniform lower sectional curvature bound organize infinitesimally into a *continuous field of locally compact metric groups* over the Alexandrov limit, generalizing the nilpotent Lie group models of Cheeger–Fukaya–Gromov without assuming an upper curvature bound?

Why is it a "good" problem:  
This problem captures a sharp conceptual gap between smooth and lower‑curvature collapsing theory: in the smooth, two‑sided curvature setting, local models are nilpotent Lie groups equipped with left‑invariant metrics; under only a lower bound, the corresponding blow‑up limits become singular Alexandrov spaces whose algebraic nature is unknown. Asking whether these limits can still be canonically described by locally compact metric groups probes the existence of an algebraic structure underlying Alexandrov collapses.  

If successful, it would unify *collapsing geometry*, *Alexandrov geometry*, and *metric group theory*, yielding a “metric N‑structure” concept for one‑sided curvature bounds. The problem is formally simple—expressed in terms of canonical local models and continuity—but it demands deep insight into how curvature bounds, geometric limits, and group‑like structures interact. It thus exemplifies a “good” research‑level problem: conceptually elegant, bridging distinct areas, and aimed at revealing new fundamental structures in Riemannian geometry.